\section*{Capítulo \ref{ch:b2:eruptive-styles} --- Las erupciones volcánicas}

\begin{solution}{exo:b2:eruptive-styles:1}
(a) $\qty{1.3}{km^3} = \qty{1.3e9}{m^3}$, entre $10^9$ y $10^{10}$: VEI~5.
(b) $\qty{2.9e8}{m^3}$: VEI~4. (c) $\qty{2e6}{m^3}$, entre $10^6$ y $10^7$:
VEI~2. (d) El VEI~6 empieza en \qty{e10}{m^3}, de modo que el volumen aparente de \omterm{def:b2:eruptive-styles:tephra}{tefra} debe ser de al menos
\qty{10}{km^3}; «más de \qty{5}{km^3}» es una cota inferior, no una contradicción.
\end{solution}

\begin{solution}{exo:b2:eruptive-styles:2}
(a) Ceniza (de menos de \qty{2}{mm}). (b) Un lapillo (\qtyrange{2}{64}{mm}).
(c) Un bloque: de más de \qty{64}{mm}, anguloso, sólido al ser expulsado. (d) Una
\omterm{def:b2:eruptive-styles:tephra}{bomba volcánica}, una bomba en huso: de más de \qty{64}{mm} y modelada mientras estaba fundida,
en vuelo.
\end{solution}

\begin{solution}{exo:b2:eruptive-styles:3}
(a) $0.304\times(10^6)^{0.241} = 0.304\times27.9 = \qty{8.5}{km}$;
$0.304\times(10^8)^{0.241} = 0.304\times84.7 = \qty{25.8}{km}$.
(b) $Q = (35/0.304)^{1/0.241} = 115.1^{4.149} = \qty{3.6e8}{kg/s}$, e $I =
8.55 + 3 = 11.6$. (c) $Q \propto H^{4.15}$, de modo que una altura un \qty{20}{\percent} demasiado alta
sobrestima $Q$ en un factor $1.2^{4.15} = 2.1$.
\end{solution}

\begin{solution}{exo:b2:eruptive-styles:4}
$V_{\mathrm{DRE}} = 0.29\times1000/2500 = \qty{0.116}{km^3}$; $m =
\num{0.29e9}\times1000 = \qty{2.9e11}{kg}$; $M = \log_{10}(\num{2.9e11}) - 7 = 11.46
- 7 = 4.46$.
\end{solution}

\begin{solution}{exo:b2:eruptive-styles:5}
$k = \ln(40/10)/(15 - 5) = 1.386/10 = \qty{0.139}{km^{-1}}$; $T_0 = 40\,e^{0.139\times5}
= 40\times2.00 = \qty{80}{cm}$; $\ln2/k = \qty{5.0}{km}$. Con $T_0 =
\qty{8.0e-4}{km}$: $V = 2\times\num{8.0e-4}/0.139^2 = \qty{0.083}{km^3} =
\qty{8.3e7}{m^3}$, VEI~3.
\end{solution}

\begin{solution}{exo:b2:eruptive-styles:6}
(a) $\int_0^x T_0e^{-ky}2y\,dy = 2T_0\big[-\tfrac{y}{k}e^{-ky} - \tfrac{1}{k^2}e^{-ky}\big]_0^x
= \tfrac{2T_0}{k^2}[1 - (1 + kx)e^{-kx}]$. (b) Iterando $u = \ln(2(1 + u))$ a partir de
$u = 1$: $1.386$, $1.564$, $1.635$, $1.662$, $1.672$, \dots, que converge a
$u = 1.678$, para el que $(1 + u)e^{-u} = 2.678\times0.1867 = 0.500$. (c) $x_{1/2} =
1.678/0.139 = \qty{12.1}{km}$, dentro de la isópaca más externa ($\sqrt A =
\qty{15}{km}$); y sin embargo el volumen más allá de esa isópaca es $(1 + 2.09)e^{-2.09} =
\qty{38}{\percent}$ del total. Casi dos quintas partes del depósito están donde
no se midió nada: el volumen descansa en la extrapolación.
\end{solution}

\begin{solution}{exo:b2:eruptive-styles:7}
(a) $p_s = (5.0/0.444)^2 = \qty{127}{MPa}$, a $127/0.025 = \qty{5.1}{km}$.
(b) $x = (5.0 - 0.444\sqrt{10})/100 = (5.0 - 1.404)/100 = 0.0360$;
$V_g = 0.0360\times8.314\times1123/(0.018015\times10^7) = \qty{1.86e-3}{m^3}$;
$V_m = 0.964/2300 = \qty{4.19e-4}{m^3}$; $\phi = 1.86/(1.86 + 0.419) = 0.82$.
(c) Un fundido basáltico es de miles a miles de millones de veces menos viscoso que uno
riolítico (\cref{ch:b2:magma-emplacement}): sus burbujas suben,
se fusionan y escapan, y el gas sale por una espuma permeable o por la
boca, de modo que el \omterm{def:b2:magma-emplacement:magma}{magma} se desgasifica en sistema abierto en lugar de conservar su gas; su
fragmentación, cuando se produce, se limita a una fuente cerca de la boca. Un ascenso más lento
deja más tiempo para ello.
\end{solution}

\begin{solution}{exo:b2:eruptive-styles:8}
(a) $u_z(0) = 0.75\times10^7/(\pi\times(5000)^2) = \qty{0.095}{m} = \qty{95}{mm}$.
(b) $z_0 = 3.8/0.766 = \qty{5.0}{km}$. (c) $(1 + 1)^{-3/2} = 0.35$ del
máximo.
\end{solution}

\begin{solution}{exo:b2:eruptive-styles:9}
El alcance máximo se da a \ang{45}: $R = v^2/g = 100^2/9.81 = \qty{1.0}{km}$, y
$200^2/9.81 = \qty{4.1}{km}$. El rozamiento del aire acorta el vuelo, de modo que son
cotas superiores, cercanas a la realidad para los bloques grandes y demasiado largas para los pequeños.
Como $R = v^2/g$, la zona de peligro crece como el cuadrado de la velocidad de lanzamiento:
duplicar la velocidad de la explosión cuadruplica su radio.
\end{solution}

\begin{solution}{exo:b2:eruptive-styles:10}
(a) Con $T = T_0e^{-k_1x}$ para $x < X$ y $T = T_Xe^{-k_2(x - X)}$ más allá, $T_X =
T_0e^{-k_1X}$ (continuidad), úsense $\int_0^X ye^{-ky}dy = \frac{1}{k^2} - e^{-kX}\big(\frac{X}{k} +
\frac{1}{k^2}\big)$ y $\int_X^\infty ye^{-ky}dy = e^{-kX}\big(\frac{X}{k} +
\frac{1}{k^2}\big)$:
la parte proximal da $2T_0\big[\frac{1}{k_1^2} - e^{-k_1X}\frac{k_1X + 1}{k_1^2}\big]$ y la
parte distal $2T_Xe^{k_2X}\int_X^\infty ye^{-k_2y}dy = 2T_X\frac{k_2X + 1}{k_2^2}$; con $T_X =
T_0e^{-k_1X}$ su suma es la fórmula del recuadro del método.
(b) $T_0 = \qty{0.01285}{km}$: $2T_0/k_1^2 = 0.0257/0.1332 = \qty{0.193}{km^3}$;
$e^{-0.365\times8.56} = 0.0440$; $(0.138\times8.56 + 1)/0.138^2 = 114.5$ y
$(0.365\times8.56 + 1)/0.365^2 = 30.96$; la corrección es $0.0257\times0.0440\times83.6
= \qty{0.094}{km^3}$ y $V = \qty{0.287}{km^3}$, los \qty{0.29}{km^3} publicados.
Más allá de $X$ quedan $2T_X(X/k_2 + 1/k_2^2) = 2\times\num{5.65e-4}\times114.5 =
\qty{0.129}{km^3}$, es decir, un \qty{45}{\percent}.
\end{solution}

\begin{solution}{exo:b2:eruptive-styles:11}
(a) $100^{1/4} = 3.16$; las leyes empíricas dan $100^{0.241} = 3.03$ y
$100^{0.226} = 2.83$. (b) $H \propto N^{-3/4}$: una atmósfera más fuertemente estratificada
detiene la columna más abajo; para una $N$ un \qty{20}{\percent} mayor,
$1.2^{-0.75} = 0.87$, una altura un \qty{13}{\percent} menor. (c) Las plumas débiles las
dobla el viento, lo que las rebaja, mientras que el aire húmedo libera calor latente en
ellas, lo que las eleva; estos efectos pesan de forma distinta en las erupciones pequeñas y en las grandes
y no están en la ley simple. Y los caudales solo se conocen con un margen de
factores de varios, un error en la variable del eje horizontal que sesga
hacia cero la pendiente ajustada.
\end{solution}

\begin{solution}{exo:b2:eruptive-styles:12}
(a) $132/400\times100 = 33$ por siglo antes de 1900; $95/125\times100 = 76$ por
siglo desde entonces. (b) $33/76 = \qty{43}{\percent}$. (c) $21/56 = \qty{38}{\percent}$:
el mismo subregistro, en un factor de dos a tres. (d) Las erupciones mayores se
ven desde más lejos y dejan depósitos más gruesos y extensos que resisten la erosión
más tiempo; una erupción de VEI~4 en una isla remota en 1700 podía pasar sin registrarse, una
de VEI~6 difícilmente.
\end{solution}

\begin{solution}{pb:b2:eruptive-styles:1}
\textbf{1.} $\ln T = 5.298$, $4.605$, $3.871$, $2.639$, $1.386$. Las pendientes entre
isópacas sucesivas son $-0.0693$, $-0.0367$, $-0.0308$ y $-0.0313$ por km: un
segmento empinado entre 10 y \qty{20}{km}, uno suave de 40 a
\qty{120}{km}; la ruptura está entre 20 y \qty{40}{km}.

\textbf{2.} $k_1 = \ln(200/100)/10 = \qty{0.0693}{km^{-1}}$; $T_0 = 200\,e^{0.693} =
\qty{400}{cm}$.

\textbf{3.} $167e^{-1.244} = \qty{48.1}{cm}$, $167e^{-2.488} = \qty{13.9}{cm}$,
$167e^{-3.732} = \qty{4.00}{cm}$: a menos de un \qty{1}{\percent} de 48, 14 y 4.0.

\textbf{4.} $400e^{-0.0693X} = 167e^{-0.0311X}$ da $X = \ln(400/167)/(0.0693 -
0.0311) = 0.873/0.0382 = \qty{22.9}{km}$.

\textbf{5.} $\ln2/0.0693 = \qty{10.0}{km}$ y $\ln2/0.0311 = \qty{22.3}{km}$.

\textbf{6.} $T_0 = \qty{0.0040}{km}$: $V = 2\times0.0040/0.0693^2 =
\qty{1.67}{km^3}$.

\textbf{7.} $e^{-0.0693\times22.9} = 0.205$; $(0.0311\times22.9 + 1)/0.0311^2 = 1770$;
$(0.0693\times22.9 + 1)/0.0693^2 = 539$; corrección $2\times0.0040\times0.205\times1231 =
\qty{2.02}{km^3}$; $V = 1.67 + 2.02 = \qty{3.68}{km^3}$, del que la corrección es
el \qty{55}{\percent}.

\textbf{8.} $3.68/1.67 = 2.2$. La ceniza distal fina es delgada, fácil de erosionar,
de removilizar, o puede caer al mar; en un depósito antiguo a menudo se ha perdido o nunca se ha
cartografiado, y el volumen queda entonces subestimado en un factor de dos o más.

\textbf{9.} $\qty{3.68e9}{m^3}$ está entre $10^9$ y $10^{10}$: VEI~5.

\textbf{10.} $V_{\mathrm{DRE}} = 3.68\times900/2400 = \qty{1.38}{km^3}$.

\textbf{11.} $m = \num{3.68e9}\times900 = \qty{3.31e12}{kg}$.

\textbf{12.} $M = \log_{10}(\num{3.31e12}) - 7 = 12.52 - 7 = 5.52$.

\textbf{13.} $3.68/1.3 = 2.8$ veces.

\textbf{14.} Una columna de más de \qty{25}{km}, y un estilo pliniano.

\textbf{15.} $Q = \num{3.31e12}/36000 = \qty{9.2e7}{kg/s}$; $I = 7.96 + 3 = 11.0$.

\textbf{16.} $0.304\times(\num{9.2e7})^{0.241} = 0.304\times83.0 = \qty{25.2}{km}$.

\textbf{17.} $0.345\times(\num{9.2e7})^{0.226} = 0.345\times63.1 = \qty{21.8}{km}$.

\textbf{18.} Los \qty{31}{km} medidos están un \qty{23}{\percent} por encima de la primera
predicción y un \qty{42}{\percent} por encima de la segunda, dentro de la
dispersión de un \qtyrange{53}{60}{\percent} de las leyes: coherente.

\textbf{19.} $Q = (31/0.304)^{1/0.241} = 102.0^{4.149} = \qty{2.2e8}{kg/s}$, 2.3 veces
el caudal medio. Los clastos registran la columna más alta, en el máximo de la
erupción, mientras que la pregunta~15 promedia sobre toda la fase; y parte de la masa
puede faltar en el depósito (ceniza distal, \omterm{def:b2:eruptive-styles:pdc}{flujos piroclásticos}).

\textbf{20.} Sí: de 22 a \qty{31}{km}, muy por encima de la tropopausa, a \qty{16}{km}.
El gas llevado a la estratosfera se extiende alrededor del planeta; el gas clave es el
dióxido de azufre, y los veinte millones de toneladas que el Pinatubo inyectó allí en 1991
enfriaron el planeta cerca de medio grado durante dos años.

\textbf{21.} $\lambda_3 = 399/75 = 5.32$, $\lambda_4 = 95/125 = 0.760$, $\lambda_5 = 22/225 =
0.0978$, $\lambda_6 = 8/525 = 0.0152$ por año; logaritmos $0.726$, $-0.119$, $-1.010$,
$-1.817$.

\textbf{22.} La recta da $10^{0.73} = 5.37$, $10^{-0.12} = 0.759$, $10^{-0.97} =
0.107$ y $10^{-1.82} = 0.0151$: a menos del \qty{1}{\percent}, el \qty{0}{\percent},
el \qty{10}{\percent} y el \qty{1}{\percent} de los recuentos.

\textbf{23.} El registro de las grandes erupciones sigue siendo completo más atrás en el tiempo,
porque se perciben desde más lejos y sus depósitos sobreviven más; una
ventana que empezara en 1500 contaría muy por debajo las erupciones de VEI~3, mientras que una ventana
que empezara en 1950 apenas contendría erupciones de VEI~6.

\textbf{24.} $\lambda_5 = \num{0.107}$ por año, un periodo de retorno de $1/0.107 =
\num{9.3}\text{ años}$; $1 - e^{-1.07} = 1 - 0.343 = \qty{66}{\percent}$ en diez años.

\textbf{25.} $\lambda_7 = 10^{3.28 - 5.95} = 10^{-2.67} = 0.00214$ por año, una cada
470 años; el registro da $2/1025 = 0.00195$ por año: coincidencia con un margen del
\qty{10}{\percent}, para un punto no usado en el ajuste.

\textbf{26.} $\lambda_6 = 10^{3.28 - 5.10} = 10^{-1.82} = 0.0151$ por año; $\lambda t = 1.51$
y $1 - e^{-1.51} = 1 - 0.220 = 0.780$: la probabilidad de al menos una erupción
de VEI~6 o más en algún lugar de la Tierra en los próximos cien años es
\textbf{\qty[text-series-to-math=true]{78}{\percent}}. Se basa en un ritmo constante,
en erupciones independientes entre sí (un proceso de Poisson) y en un
catálogo completo dentro de sus ventanas.
\end{solution}
